\documentstyle[%


righttag%


,11pt%


,amssymb%


,russian%               remove in Eng.


,izvru-ks%                 Set izven% in Eng.


]{amsart}





% gil@comtv.ru





\newcommand{\Ann}{\operatorname{Ann}}


\newcommand{\tr}{\operatorname{tr}}


\newcommand{\reg}{\operatorname{reg}}


\newcommand{\tts}[1]{}





\theoremstyle{plain}


\theorembodyfont{\it}


\newtheorem{asrtnonum}{\hskip\parindent Утверждение}


\renewcommand{\theasrtnonum}{}





\theoremstyle{definition}


\theorembodyfont{\rm}


\newtheorem{notenonum}{\hskip\parindent Замечание}


\renewcommand{\thenotenonum}{}





%\hoffset=-15mm%                  For Eng.


\setcounter{page}{76}


\god{2005}


\nomer{6~(517)} %                        Remove in Eng.


%\nomer{Vol.\,49, No.\,6}%               Set in Eng.


%\str{\pageref{begin}--\pageref{end}}%  Set in Eng.


\udk{514.144}





\author{\it Ю.Л.\,Гилуч}


% NB:   ^^ remove in Eng.


\title{Вещественный аналог преобразования Брайанта\\ и рациональные интегральные 


кривые\\ заданного распределения в $\Bbb{P}^3$}





\date{}


%\pagestyle{myheadings}%         Set in Eng.


\pagestyle{plain} %              Remove in Eng.


\begin{document}


%\thispagestyle{plain}%          Set in Eng.


\maketitle


\label{begin}











{\bf 1. Введение.}


Особенностью проведенного в данной работе исследования контактного распределения 


на $\Bbb{P}^3$ является сопоставление геометрии 


этого распределения с геометрией вложенного в $\Bbb{P}^3$ интегрального 


многообразия. 





Наиболее интересными однородными трехмерными вещественными контактными 


многообразиями являются 1)~$(\Bbb{P}(T^{\ast}\Bbb{P}^2),\theta)$,


где $\theta=x^0dy_0+x^1dy_1+x^2dy_2$; 2)~$(\Bbb{P}^3,\omega)$,


где $\omega=x^3dx^0-3x^2dx^1+3x^1dx^2-x^0dx^3$.





По теореме Дарбу ([1], с.\,328) трехмерные контактные 


многообразия локально изоморфны. 


Поскольку $(\Bbb{P}^3,\omega)$ и $(\Bbb{P}(T^{\ast}\Bbb{P}^2),\theta)$ 


являются примерами алгебраических многообразий, естественно поставить вопрос о 


существовании бирационального изоморфизма между ними,


т.\,е. такого не всюду определенного изоморфизма, который сохраняет 


контактные структуры и задается рациональными 


функциями. Простой бирациональный изоморфизм, сохраняющий контактные 


структуры, был предложен Робертом Брайантом в [2], где решался 


вопрос о конформном представлении двумерных минимальных поверхностей и доказано, 


что все они получаются из римановых поверхностей. Важным техническим приемом был 


построенный там же бирациональный изоморфизм $f$ из $\Bbb{P}(T^{\ast}\Bbb{C}


\Bbb{P}^2)$ в $\Bbb{C}\Bbb{P}^3$ со стандартными контактными 


формами и доказанная им








\begin{thm}[{\series{m}\shape{n}\selectfont [2]}] 


Пусть $C$ --- контактная кривая в $\Bbb{CP}^3$. 


Тогда $C$ является прямой или имеет вид $f(\wt{D})$, где  $\wt{D}\subset\Bbb{P}


(T^{\ast}\Bbb{C}\Bbb{P}^2)$ --- горизонтальное поднятие приведенной и 


неприводимой плоской кривой $D\subset\Bbb{CP}^2$  степени, по крайней мере, $2$.


\end{thm}





С помощью изоморфизма $f$ каждая комплексная рациональная интегральная кривая 


получалась из рациональной алгебраической кривой на комплексной проективной 


плоскости с помощью двух операций: 1)~поднятие кривой в $\Bbb{P}(T^{\ast}


\Bbb{P}^2)$; 2)~бирациональное отображение этого многообразия в $\Bbb{P}^3$.





Метод Брайанта позволяет решать различные вопросы. Так, проблемой является 


построение явных примеров однородных гладких лежандровых многообразий, которая 


легко разрешима [3] с применением конструкции Брайанта.





Основной целью данной работы является перенесение конструкции Брайанта в 


вещественную область и ее геометрическая интерпретация. В п.\,2 найден 


вещественный аналог преобразования Брайанта.





Начиная с работ А.\,Фосса, Д.Н.\,Синцова [4], исследование контактных структур 


на трехмерных многообразиях проводилось с помощью изучения интегральных кривых 


заданного распределения. Поскольку контактная структура в $\Bbb{P}^3$ 


определяется нормкривой, показалось естественным охарактеризовать интегральные 


кривые и их частные классы, используя непосредственно нормкривую. Эта задача 


решается в п.\,3.





Бирациональное преобразование, предложенное Брайантом, связывает геометрию обоих 


многообразий и сводит в ряде случаев вопросы, касающиеся интегральных кривых, к 


вопросам о кривых на проективной плоскости. Этот факт позволяет решать задачу о 


соединимости двух точек трехмерного проективного пространства рациональной 


интегральной кривой. Вопрос о соединимости двух точек кусочно интегральной кривой 


поставлен и решен Чжоу (1939) и Рашевским (1938). В 1963~году Смэйл (Smale) 


поставил вопрос о возможности соединить две точки $C^{\infty}$-гладкой кривой. 


Эта задача была решена В.Н.\,Черненко [5]. В этой задаче идет речь о 


произвольной контактной структуре на произвольном гладком многообразии. 


Если в $\Bbb{P}^3$ задана контактная структура, то естественно поставить 


вопрос о соединимости двух точек рациональной интегральной кривой и определить 


наименьшую степень кривой, достаточной для решения 


задачи. Эта задача решается в п.\,4.


\medskip








{\bf 2. Вещественный аналог преобразования Брайанта.}


Пусть $\Bbb{CP}^3$ --- трехмерное комплексное проективное пространство. 


На нем задана стандартная структура, которая в аффинной карте $(1:z_1:z_2:z_3)$


имеет вид $\theta=dz_1-z_3dz_2+z_2dz_3$. Пусть $D\subset\Bbb{CP}^2$ --- 


приведенная и неприводимая над $\Bbb C$ кривая, и $D_{\reg}$ --- гладкая часть $D$. 


Тогда точки $D_{\reg}$ вместе с касательными направлениями формируют кривую в 


$\Bbb{P}(T^{\ast}\Bbb{P}^2)$. Это дает поднятие $D_{\reg}$ в $\Bbb{P}


(T^{\ast}\Bbb{P}^2)$. Замыкание Зариского этого поднятия в $\Bbb{P}


(T^{\ast}\Bbb{P}^2)$ есть поднятие $D$ в $\Bbb{P}(T^{\ast}\Bbb{P}^2)$. 


Оно называется [6] горизонтальным поднятием $D$ и обозначается $\wt


{D}$. Кривая $\wt{D}$ гладкая, если кривая $D$ имеет только неразветвленные или


простые особенности типа клюва.





Брайант определил [2] бирациональное контактное отображение $f:\Bbb{P}


(T^{\ast}\Bbb{P}^2)\to\Bbb{CP}^3$, полагая $(x:y,[\lambda_1:\lambda_2])$ 


как $(\lambda_1:x\lambda_1-\frac{1}{2}y\lambda_2:y\lambda_1:\frac{1}{2}\lambda_2


)$, где $(x,y)$ --- координаты на $\Bbb{C}^2\subset\Bbb{CP}^2$, а $(\lambda_


1,\lambda_2)$ --- однородные координаты слоя.





Существует естественная контактная структура на $\Bbb{P}(T^{\ast}\Bbb


{P}^2)$ [6]. Бирациональное отображение $f$ контактно в том смысле, что оно 


отображает контактные кривые в контактные кривые. Пусть $U_2=\{x_2=0\}\subset


\Bbb{CP}^2$ --- аффинное открытое множество с аффинными координатами $(x,y)=


(x_0/x_2,x_1/x_2)$. На $U_2$ мы можем отождествить слоевые координаты $[\lambda


_1,\lambda_2]$


с


$[-y_0,y_1]$. Таким образом, преобразование Брайанта может быть


переопределено в виде [6]


$$


f([x_0: x_1: x_2],[y_0: y_1: y_2])=[2x_2y_0:2x_0y_0


+x_1y_1:2x_1y_0:-x_2y_1].


$$





\begin{thm}[{\series{m}\shape{n}\selectfont [6]}]


Пусть $C$ --- контактная кривая в $\Bbb{CP}^3$. Тогда $C$ 


является прямой или имеет вид $f(\wt{D})$, где  $\wt{D}\subset\Bbb{CP}^2$ --- 


горизонтальное поднятие приведенной и неприводимой плоской кривой $D$ степени,


по крайней мере, $2$.


\end{thm}





Теорема Брайанта сводит исследование контактных кривых в $\Bbb{CP}^3$ к 


исследованию алгебраических кривых в $\Bbb{CP}^2$ и их поднятий в $\Bbb{P}


(T^{\ast}\Bbb{P}^2)$. Важным свойством отображения $f$ является то, что


оно является отображением, сохраняющим контактные формы, и это соответствие 


является рациональным. 





Перенесем преобразование Брайанта 


в вещественную область и дадим его геометрическую характеристику.





Пусть $\Bbb{P}^3=\Bbb{P}(\Bbb{R}^4)$ --- трехмерное вещественное 


проективное пространство, порожденное векторным пространством $\Bbb{R}^4$.





Кокасательное пространство $T^{\ast}_A\Bbb{RP}^3$ к проективному пространству 


$\Bbb{RP}^3$ в точке $A=p(a)$ будем отождествлять с подпространством $\Ann 


\langle a\rangle\subset(\Bbb{R}^4)^{\ast}$, т.\,е. $T^{\ast}_A\Bbb{RP}^3$ 


будем рассматривать как совокупность линейных форм на $\Bbb{R}^4$, равных нулю 


на векторе $a$.





Проективизация кокасательного расслоения  


$$


\Bbb{P}(T^{\ast}\Bbb{P}^2)=\bigcup\limits_{A\in\Bbb{P}^2}\Bbb{P}(T^{\ast}_A


\Bbb{P}^2)


$$


--- многообразие всех проективных пространств $\Bbb{P}(T^{\ast}_A


\Bbb{P}^2)$, порожденных кокасательными пространствами $T^{\ast}_A\Bbb{P}^2$ 


к $\Bbb{P}^2$ во всех его точках $A\in\Bbb{P}^2$.








\begin{defn}[{\series{m}\shape{n}\selectfont [7], с.\,289}]


{\it Нуль-парой} в $\Bbb{P}^2$ называется  


пара $(A,\Xi)$, где $A\in\Bbb{P}^2$ и $\Xi\in\check{\Bbb{P}}^2$. Нуль-пара 


называется {\it вырожденной}, если точка $A=p(a)$ и прямая $\Xi=p(\xi)$ 


инцидентны, т.\,е. $A\in\Xi$ (или, что эквивалентно, $\xi(a)=0$).


\end{defn}





Известно [8], что {\it многообразие вырожденных нуль-пар диффеоморфно 


проективизации $\Bbb{P}(T^{\ast}\Bbb{P}^2)$ кокасательного расслоения 


пространства $\Bbb{P}^2$}.





На многообразии $\Bbb{P}(T^{\ast}\Bbb{P}^2)$ каноническая структура 


задается уравнением


$$


\theta=x^0dy_0+x^1dy_1+x^2dy_2.


$$


  


Многообразие $X_3=\Bbb{P}(T^{\ast}\Bbb{P}^2)$ можно 


задать в произведении $\Bbb{P}^2\times\check{\Bbb{P}}^2$ биоднородным уравнением 


$x^iy_i=0$, где $M=(x^0: x^1: x^2)\in\Bbb{P}^2$, $l=(y_0


: y_1: y_2)\in\check{\Bbb{P}}^2$, а $\check{\Bbb{P}}^2$


--- проективизация $(\Bbb{R}^3)^{\ast}$.





Для того чтобы представить произведение проективных алгебраических 


многообразий в виде проективного пространства, используем вложение Сегре


([9], с.\,74) в инвариантной форме $s:\Bbb


{P}^2\times\check{\Bbb{P}}^2\rightarrow\Bbb{P}(\frak{gl}(3,\Bbb


{R}))$. Если $(M,l)\in\Bbb{P}^2\times\check{\Bbb{P}}^2$, полагаем 


$$


s(M,l)=[x^iy_je_i\otimes e^j]\in\Bbb{P}^8,


$$


где $e_0$, $e_1$, $e_2$ и $e^0$, $e^1$, $e^2$ --- двойственные 


базисы $\Bbb{R}^3$ и $(\Bbb{R}^3)^{\ast}$, а $\Bbb


{P}^8$ --- проективизация алгебры Ли $\frak{gl}(3,\Bbb{R})$ вещественных


$(3\times 3)$-матриц. 





Уравнение $\tr M=0$, $M\in\frak{gl}(3,\Bbb{R})$, задает в $\Bbb{P}^8$


гиперплоскость $\Bbb{P}^7$, и ограничение $s$ на $X_3$ определяет вложение $X_


3$ в проективизацию алгебры Ли $\frak{sl}(3,\Bbb{R})$ \ $s:X_3\rightarrow


\Bbb{P}^7=\Bbb{P}(\frak{sl}(3,\Bbb{R}))$. Это вложение не зависит 


от выбора проективного репера в $\Bbb{P}^2$.








\begin{thm}


Образ $s$ в $\frak{sl}(3,\Bbb{R})$ состоит из классов 


эквивалентности матриц A таких, что $\bigwedge\limits^2 A=0$, $\tr A=0$,


$A\neq 0$.


\end{thm}








\begin{thm}


Проективное касательное пространство $\Bbb{P}(T_MX_3)$


в точке $M\in s(X_3)$ изоморфно проективизации пространства всех квадратных $(3


\times 3)$-матриц A с нулевым следом, удовлетворяющих условию $AM+MA=0$.


\end{thm}





Будем называть векторное пространство $T=\{A\mid AM+MA=0$, $\tr A=0\}$


{\it аналитическим касательным пространством} к многообразию $X_3$ в точке $M$,


поскольку по теореме 4 его проективизация совпадает с проективизацией 


касательного пространства.





Пусть $\Bbb{P}(T^{\ast}\Bbb{P}^2)@>{s\;}>>\Bbb{P}^7$ --- 


вложение Сегре и 


$$


s(x^0:x^1:x^2;y_0:y_1:y_2)=\begin{pmatrix}


-x^1y_1-x^2y_2 & x^0y_1 & x^0y_2 \\


x^1y_0 & x^1y_1 & x^1y_2 \\


x^2y_0 & x^2y_1 & x^2y_2


\end{pmatrix}.


$$





Поскольку на $\frak{sl}(3,\Bbb{R})$ имеется каноническая квадратичная


форма (форма Киллинга), можно с ее помощью задать


кососимметрическую билинейную форму (форму Кириллова--Сурио) [1]: если 


$A\in X_3$, то кососимметрическая форма задается из условия


$$


\Omega(U,V)=B(A,[U,V])=6\tr(A\cdot[U,V]),\ \ U,V\in X_3. 


$$





Ограничение этой кососимметрической формы на касательное пространство 


$\Bbb{P}^3=\Bbb{P}(T_{M}Y)$, $M\in s(X_3)=Y$, позволяет задать в 


проективизации касательного пространства $\Bbb{P}(T_MY)$ структуру 


контактного многообразия.





Сформулируем основной результат.








\begin{thm}


Существуют точка $M\in s(X_3)=Y$ и проекция $\tau$ на касательное 


пространство $\Bbb{P}(\frak{sl}(3,\Bbb{R}))\rightarrow\Bbb{P}(T_MY)=


\Bbb{P}^3$ такие, что преобразование $s\circ\tau$ переводит всякую 


алгебраическую кривую на $\Bbb{P}^2$ в интегральную кривую 


в пространстве $\Bbb{P}^3$, 


снабженном канонической контактной структурой $\Omega$.


\end{thm}








{\bf 3. О рациональных интегральных кривых уравнения Пфаффа в $\Bbb{P}^3$.}


Исследуются рациональные интегральные кривые уравнения 


$\omega=x^3dx^0-3x^2dx^1+3x^1dx^2-x^0dx^3$,


заданного на вещественном проективном пространстве $\Bbb{P}^3$. 








\begin{defn}


{\it Нормкривой} $C^3$ будем называть рациональную 


алгебраическую кривую третьей степени, которая в некотором проективном репере 


представима в виде $x^0:x^1:x^2:x^3=1:t:t^2:t^3$.


\end{defn}





Известно ([10], с.\,193), что нормкривая в проективном пространстве $\Bbb{P}^3$ 


определяет корреляцию, при которой произвольной 


точке $X=(x^0:x^1:x^2:x^3)\in\Bbb{P}^3$ отвечает гиперплоскость с уравнением


$$


\sum_{i=0}^{3}a_{i}y^{i}=0,\ \ \text{где} \ \


a_{i}=(-1)^{i}


\begin{pmatrix}3\\i\end{pmatrix}


x^{3-i},\quad i=0,1,2,3,


$$


т.\,е. уравнение имеет вид


$$


\Pi(X): x^3y^0-3x^2y^1+3x^1y^2-x^0y^3=0.


$$








\begin{thm}


Проективизация контактной плоскости уравнения $\omega=0$


совпадает в каждой точке $X\in\Bbb{P}^3$ c плоскостью корреляции $\Pi(X)$, 


задаваемой нормкривой $C^3$.


\end{thm}





Эта теорема позволяет следующим образом охарактеризовать интегральные кривые 


в $\Bbb{P}^3$.








\begin{thm}


Пусть задана некоторая кривая $\gamma(t)$, $l(t)$ --- 


касательная в точке $M\in\gamma(t)$ и пучок плоскостей $\pi(t,\lambda)$, 


осью которого является прямая $l(t)$. Если для произвольного $t$ найдется такое  


$\lambda=\lambda(t)$, что соприкасающиеся плоскости к нормкривой в точках 


пересечения $\pi(t,\lambda(t))$ c $C^3(t)$ пересекаются в точке $M$, то 


$\gamma(t)$ является интегральной кривой уравнения $\omega=0$. Верно и обратное.


\end{thm}








\begin{asrtnonum}


Для того чтобы точки пересечения плоскости 


корреляции с нормкривой $C^3(t)$ являлись вещественными, необходимо и 


достаточно, чтобы интегральная кривая располагалась в области


$$


3(x^1)^2(x^2)^2-4x^0(x^2)^3+6x^0x^1x^2x^3-4(x^1)^3x^3-(x^0)^2(x^3)^2\leq 0.


$$


\end{asrtnonum}





Рациональная интегральная кривая (здесь и далее, если не отмечено специально, 


кривая является интегральной для уравнения $\omega=0$) не является 


произвольной и в топологическом смысле.  





Рассмотрим три точки, принадлежащие нормкривой $C^3(t)$,


$$


M_1=(1:t_1:t_1^2:t_1^3),M_2=(1:t_2:t_2^2:t_2^3),M_3=(1:t_3:t_3^2:t_3^3).


$$








\begin{thm}


Три рациональные функции $t_1(t),t_2(t),t_3(t)$ определяют 


интегральную кривую тогда и только тогда, когда выполняется условие


$$


t'_1(t_3-t_2)^2+t'_2(t_1-t_3)^2+t'_3(t_2-t_1)^2=0.


$$


Здесь $t'_i=dt_i/dt$.


\end{thm}





Заметим, что, задав $t_2$, $t_3$, получим дифференциальное уравнение, 


разрешенное относительно производной, и, следовательно, имеющее локальное 


решение. Для того чтобы решение было глобальным, должны быть сделаны 


некоторые ограничения. 





Сформулируем способ построения интегральной кривой.








\begin{thm} 


Если $v$, $w$ --- рациональные функции и выражение 


$$


V=\frac{v'w^2-w'v^2}{v^2+w^2+(v+w)^2}


$$ 


не имеет простых полюсов, то формулами 


$$


t_1\ \; \text{--- первообразная от $V$},\quad t_2=t_1+w,\ \ t_3=t_1-v


$$


задается рациональная интегральная кривая следующего вида\/${:}$


$$


x^0:x^1:x^2:x^3=3:t_1+t_2+t_3:t_1t_2+t_1t_3+t_2t_3:3t_1t_2t_3.


$$


\end{thm}





\begin{notenonum}


Заметим, что, полагая в уравнении 


$2t'_1(v^2+w^2+vw)=v'w^2-w'v^2$,


$\tau=\frac{vw}{v-w}$, $t_1(t)=t$, приходим к 


уравнению $d\tau/dt=-2-6v\tau-6v^2\tau^2$,


которое является уравнением Риккати


$dx/dt=P(t)+Q(t)x+R(t)x^2$,


где $x=\tau$, $P(t)=-2$, $Q(t)=-6v(t)$, $R(t)=-6v^2(t)$.


\end{notenonum}





Уравнение Риккати --- это своеобразная реализация 


группы проективных преобразований. 


Отражением этого факта является теорема о постоянстве ангармонического отношения 


четырех решений уравнения Риккати [11]. С помощью этой теоремы 


можно утверждать следующее.








\begin{thm}


Если три интегральные кривые формы $\omega$ в проективном 


пространстве $\Bbb{P}^3$ определяются тремя решениями уравнения Риккати  


$$


\frac{d\tau}{dt}=-2-6v(t)\tau-6v^2(t)\tau^2,


$$


то можно построить четвертую интегральную кривую.


\end{thm}





Нормкривая $C^3$ позволяет связать с интегральной кривой проективный инвариант 


в ``первой дифференциальной окрестности''.





Рассмотрим произвольную интегральную кривую $\gamma(t)$. Тогда в каждой плоскости 


корреляции, связанной с текущей точкой $M(t)$ кривой, определены четыре прямые. 


Первые три прямые $MM_1$, $MM_2$, $MM_3$ соответствуют 


$M_1$, $M_2$, $M_3$ --- 


точкам пересечения плоскости корреляции с нормкривой $C_3$, а четвертая прямая 


$MM_4$ является касательной к $\gamma(t)$ в $M(t)$ ($M_4$ --- произвольная точка 


касательной).





Пусть $M=(m^0:m^1:m^2:m^3)$, $M_i=(m^0_i:m^1_i:m^2_i:m^3_i)$, $i=1,2,3,4$. Точке 


$M$ соответствует плоскость корреляции $\Pi(M)$. Пусть координаты плоскости 


$\Pi(M)$ есть $a=(a_0:a_1:a_2:a_3)$.








\begin{thm}


Сложное отношение четырех  инвариантных прямых в плоскости корреляции, 


связанных с текущей точкой кривой, вычисляется по формуле


$$


D=(MM_1,MM_2;MM_3,MM_4)=\frac{\det(M,a,M_1,M_3)}


{\det(M,a,M_1,M_4)}\cdot


\frac{\det(M,a,M_2,M_3)}


{\det(M,a,M_2,M_4)},


$$


где $\det(M,a,M_i,M_j)$, $i,j=1,2,3,4$, есть определитель, в строках которого 


стоят координаты точек $M$, $a$, $M_i$, $M_j$.


\end{thm}








{\bf 4. Cоединимость точек проективного пространства рациональной интегральной 


кривой.} 


В классе рациональных кривых можно решать задачу о соединимости двух точек 


трехмерного проективного пространства $\Bbb{P}^3$ рациональной интегральной 


кривой заданного распределения.





{\it Задача.} Пусть в $\Bbb{P}^3$ заданы две точки $M_1$, $M_2$. Найти 


кривую $\gamma$, удовлетворяющую следующим условиям:


\begin{itemize}


\item[(1)] $M_1,M_2\in\gamma$;


\item[(2)] $\gamma$ яляется интегральной кривой заданного распределения;


\item[(3)] $\gamma$ касается в точках $M_1$ и $M_2$ прямых $L_1$ и $L_2$, лежащих 


в плоскостях корреляции точек $M_1$, $M_2$ соответственно.


\end{itemize}





Если опустить условие (3), тогда через две любые точки в $\Bbb{P}^3$ 


проходит рациональная алгебраическая кривая степени 4, удовлетворяющая условию 


(2), --- нормкривая $C^3$. Такие кривые образуют однопараметрическое семейство. 


Степень рациональной кривой, решающей общую задачу, равна 9.





Прямая в $\Bbb{P}^3$ с заданной на ней точкой 


преобразуется в кривую на $X_3$, являющуюся поднятием некоторой кривой из 


$\Bbb{P}^2$. Касание первого порядка в $\Bbb{P}^3$ с прямой 


интерпретируется на $\Bbb{P}^2$ как касание второго порядка с кривой на 


плоскости в заданной точке.   








\begin{thm}


Всякие две точки проективного пространства $\Bbb{P}^3$ 


можно соединить рациональной интегральной кривой, удовлетворяющей условию $(3)$, 


степени меньше либо равной $9$. 


\end{thm}


 





Результаты статьи были доложены автором на трех математических конференциях


[12]--[14].








\begin{thebibliography}{99}





\bibitem{2} %1


Арнольд В.И. {\it Математические методы классической механики}. -- 


М.: Наука, 1974. -- T.\,2. -- 432~с.


\bibitem{3} %2


Bryant R. {\it Conformal and minimal immersions of compact surfaces 


into the four sphere} // J. Different. Geom. -- 1982. -- V.\,17. --


P.\,455--473. 


\bibitem{5} %3


Landsberg J.M., Manivel L. {\it Legendrian varieties}.


http://ru.arxiv.org/alg-geom/0407279.pdf.


\bibitem{6} %4


Синцов Д.М. {\it Работы по неголономной геометрии}. -- Киев: Вища школа,


1972. -- 391~с.


\bibitem{7} %5


Черненко В.Н. {\it О соединимости двух точек трехмерного 


$C^{\infty}$-многообразия интегральной кривой не вполне


интегрируемой дифференциальной 


системы} // Геометрич. сб. -- Томск: Изд-во Томск. ун-та. -- 1978. --


Т.\,19. -- C.\,3--11.


\bibitem{4} %6


Yun-Gang Ye. {\it Complex contact threefolds and their contact curves}.


http://ru.arxiv.org/alg-geom/9202003.pdf.


\bibitem{8} %7


Розенфельд Б.А. {\it Многомерные пространства}. -- М.: Наука, 1966. --


530~с.


\bibitem{9} %8


Коннов В.В. {\it Касательное расслоение проективного пространства и 


многообразие невырожденных нуль-пар} // Матем. заметки. -- 2001. --


Т.\,70. -- \No\,5. -- С.\,718--735.


\bibitem{10} %9


Шафаревич И.Р. {\it Основы алгебраической геометрии}. T.\,1. -- М.: Наука, 1988.


-- 352~с.


\bibitem{11} %10


Burau W. {\it Mehrdimensionale projektive und h\"ohere 


Geometrie}. -- Berlin, 1961. -- 436~S.


\bibitem{12} %11


Ибрагимов Н.Х. {\it Групповой анализ обыкновенных дифференциальных уравнений 


и принцип инвариантности в математической физике} // УМН.


-- 1992. -- Т.\,47. -- Вып.\,4. -- C.\,83--144.


\bibitem{13} %12


Гилуч Ю.Л. {\it Вещественный аналог преобразования Брайанта} // Материалы 


международн. молодежной научной школы-конф. ``Лобачевские чтения-2002'' --


Казань, 2002. -- C.\,20--21.


\bibitem{14} %13


Гилуч Ю.Л. {\it Неголономная вещественная алгебраическая поверхность 


в $\Bbb{RP}^3$} // Тез. докл. 5-й международной конф. по геометрии и 


топологии. -- Черкассы (Украина), 2003. -- C.\,26--28.


\bibitem{15} %14


Гилуч Ю.Л. {\it О рациональных интегральных кривых уравнения Пфаффа в 


трехмерном проективном пространстве} // Материалы третьей всероссийск. молодежной 


научн. школы-конф. ``Лобачевские чтения-2003'' -- Казань, 2003.


-- C.\,96--98.





\end{thebibliography}





\vskip 2cm





\post{\it Томский государственный}{\it Поступила}


\post{\it университет}{05.10.2004}





\label{end}


\end{document}





